% ============================================================
%  chapters/preface.tex
% ============================================================

\chapter*{Preface}
\addcontentsline{toc}{chapter}{Preface}
\markboth{Preface}{Preface}

This work begins with a suspicion: that arithmetic, gesture, language,
cognition, computation, and cosmology are not separate domains connected by
analogy, but \emph{different coordinate systems for the same underlying
geometry}.

The geometry in question is not Riemannian in the standard sense. It is the
geometry of \emph{accessibility} — of what states are reachable from here,
under what constraints, along which trajectories, with how much remaining
freedom. Five frameworks developed independently across a decade of work
converge, when viewed at the right angle, on this single organizing principle:

\medskip
\begin{tcolorbox}[
  enhanced, colback=RSVPdeep, colframe=RSVPdeep,
  coltext=white, arc=6pt, boxrule=0pt,
  left=8mm, right=8mm, top=6mm, bottom=6mm
]
  \begin{center}
    \large
    \textbf{Spherepop} (collapse) \quad
    \textbf{Gesture} (trajectory) \quad
    \textbf{Admissibility} (constraint)\\[4pt]
    \textbf{RSVP} (accessibility) \quad
    \textbf{Semantic Infrastructure} (composition)
  \end{center}
  \medskip
  \normalsize\itshape
  These are not five theories. They are five projections of one theory onto
  five different observable screens.
\end{tcolorbox}
\medskip

Part~I of this monograph builds the foundations from scratch, requiring only
mathematical maturity and patience. Later parts assume familiarity with the
earlier material. The appendices collect the formal machinery for readers who
prefer to start from definitions.

\medskip
\noindent\textit{A note on style.} The prose aims to be precise without being
arid. Boxed definitions, theorems, and examples are formally numbered. Asides
and historical notes are set off visually but are not peripheral — they are
often where the conceptual work happens. Exercises are provided at the end of
most chapters; starred exercises are open problems.

\bigskip
\begin{flushright}
  \textit{Flyxion}\\
  \textit{Canada, \the\year}
\end{flushright}
